\documentclass[11pt, letterpaper]{article}
\usepackage{mobi-lectura}

\title{Aplicación del Teorema del Límite Central (TLC)\\
Promedio de costos de importación a través de múltiples envíos}
\author{Alejandra Tabares Pozos}
\date{}

\begin{document}

\maketitle
\tableofcontents
\newpage

\section{Contexto para la Lectura}
Para ilustrar el Teorema del Límite Central (TLC), debemos cambiar la perspectiva:
ya no nos interesa solo el costo de un envío (cuya distribución es asimétrica y compleja),
sino el comportamiento del promedio de costos a lo largo de múltiples envíos.

En la sección anterior, vimos que la distribución de costos de un solo envío tenía una cola larga (asimetría positiva).
Sin embargo, la gerencia financiera no evalúa el desempeño basándose en un solo envío, sino en el presupuesto anual o trimestral.
El Teorema del Límite Central nos asegura que, aunque el costo individual no sea Normal, el promedio de costos de $n$ envíos
sí se comportará como una distribución Normal (para $n$ suficientemente grande, usualmente $>30$).

\section{Fundamento Teórico breve: por qué el TLC aplica aquí}
Modelemos el costo de un envío como una variable aleatoria $X$ (no necesariamente Normal), con media y varianza finitas:
\[
\E[X]=\mu,\qquad \Var(X)=\sigma^2<\infty.
\]
Consideremos $X_1,\dots,X_n$ i.i.d.\ como costos de $n$ envíos, y su promedio:
\[
\bar X_n=\frac{1}{n}\sum_{i=1}^n X_i.
\]
El Teorema del Límite Central establece que, al estandarizar el promedio,
\[
\frac{\sqrt{n}\,(\bar X_n-\mu)}{\sigma}\;\xrightarrow{d}\;\N(0,1),
\qquad \text{cuando } n\to\infty.
\]
En términos prácticos, para $n$ moderadamente grande, esto se usa como aproximación:
\[
\bar X_n \approx \N\!\left(\mu,\frac{\sigma^2}{n}\right).
\]
\paragraph{Cuando $\sigma$ es desconocida.}
En aplicaciones reales se reemplaza $\sigma$ por la desviación estándar muestral $s$; para $n$ grande, la aproximación Normal sigue siendo útil,
y de manera clásica se usa la corrección con distribución $t$ de Student:
\[
T=\frac{\bar X_n-\mu}{s/\sqrt{n}} \approx t_{n-1}.
\]
Esto habilita intervalos de confianza y pruebas de hipótesis sobre $\mu$ incluso cuando la distribución de $X$ es sesgada, siempre que $n$ sea suficientemente grande
y no haya varianza infinita.

\section{Datos de la Muestra}

El auditor extrajo una muestra aleatoria de $n = 40$ envíos del registro histórico de los últimos tres años. Los costos (en millones de COP) son:

\begin{center}
\small
\renewcommand{\arraystretch}{1.15}
\begin{tabular}{cccccccccc}
\toprule
\multicolumn{10}{c}{\textbf{Costo por envío (millones COP)}} \\
\midrule
31.4 & 24.2 & 29.5 & 22.8 & 35.1 & 26.9 & 28.6 & 33.0 & 21.5 & 27.4 \\
30.8 & 25.3 & 32.7 & 23.0 & 28.2 & 27.1 & 31.6 & 23.4 & 29.1 & 24.7 \\
26.8 & 34.5 & 23.1 & 30.2 & 25.6 & 28.0 & 33.4 & 21.0 & 29.9 & 26.1 \\
24.4 & 31.2 & 27.7 & 34.0 & 23.5 & 28.4 & 25.0 & 30.5 & 26.3 & 25.2 \\
\bottomrule
\end{tabular}
\end{center}

Los estadísticos de la muestra son:
\[
n = 40, \qquad \bar{x} = 27{,}78 \text{ millones COP}, \qquad s = 3{,}76 \text{ millones COP}.
\]

\section{Preguntas para Aplicación del TLC}

\subsection*{Pregunta 1: Estimación por Intervalo (Auditoría)}
Un auditor interno selecciona una muestra aleatoria de 40 envíos históricos y encuentra un costo promedio muestral ($\bar{x}$) y una desviación estándar muestral ($s$).

\textbf{Requerimiento:} Construya un intervalo de confianza del 95\% para el verdadero costo promedio de importación.
¿Está el presupuesto de $29$ millones dentro de este intervalo?

\paragraph{Nota técnica (fórmula):}
Un IC del 95\% para $\mu$ (usando $t$) es:
\[
\bar x \pm t_{0.975,\;39}\,\frac{s}{\sqrt{40}}.
\]
El presupuesto $29{,}000{,}000$ (COP) está ``dentro'' si cae entre los límites inferior y superior del intervalo.

\subsection*{Pregunta 2: Prueba de Hipótesis (Evaluación de Desempeño)}
El nuevo Director de Logística afirma que sus optimizaciones han logrado reducir el costo promedio real a menos de $29{,}000{,}000$.

\textbf{Requerimiento:} Utilizando la misma muestra de 40 envíos, realice una prueba de hipótesis con un nivel de significancia $\alpha = 0.05$.
($H_0: \mu \ge 29M$ vs $H_1: \mu < 29M$).

\paragraph{Nota técnica (estadístico y regla):}
\[
T=\frac{\bar x-\mu_0}{s/\sqrt{40}},\qquad \mu_0=29{,}000{,}000.
\]
Rechazar $H_0$ a nivel $\alpha$ si $T < t_{\alpha,\;39}$ (cola izquierda), o equivalentemente si el $p$-valor es $<\alpha$.

\subsection*{Pregunta 3: Diseño de Experimentos (Tamaño de Muestra)}
Para el próximo año fiscal, la gerencia desea estimar el costo promedio con un error máximo de $\pm \$500{,}000$ COP y una confianza del 95\%.

\textbf{Requerimiento:} ¿Cuántos envíos mínimos ($n$) deben monitorearse para garantizar esta precisión?

\paragraph{Nota técnica (diseño):}
Si se requiere un margen de error $E$ para un IC del 95\% y se aproxima con Normal,
\[
E \approx z_{0.975}\,\frac{\sigma}{\sqrt{n}}
\quad\Rightarrow\quad
n \approx \left(\frac{z_{0.975}\,\sigma}{E}\right)^2,
\]
con $z_{0.975}\approx 1.96$. En la práctica se usa una estimación piloto de $\sigma$ (por ejemplo $s$ histórico) y se redondea hacia arriba.

\section{Respuestas de Referencia}

\begin{resultadobox}[Pregunta 1: Intervalo de Confianza del 95\%]
Con $\bar{x} = 27{,}78$, $s = 3{,}76$, $n = 40$ y $t_{0.975,\,39} \approx 2{,}023$:
\[
SE = \frac{3{,}76}{\sqrt{40}} = \frac{3{,}76}{6{,}325} \approx 0{,}595 \text{ M COP}
\]
\[
IC_{95\%} = 27{,}78 \pm 2{,}023 \times 0{,}595 = 27{,}78 \pm 1{,}20 = [26{,}58\text{ M},\; 28{,}98\text{ M}]
\]
El presupuesto de 29 millones \textbf{queda fuera} del IC (el límite superior es 28.98 M). Hay evidencia estadística de que el costo promedio real es inferior a 29 M.
\end{resultadobox}

\begin{resultadobox}[Pregunta 2: Prueba de Hipótesis ($H_0\!: \mu \ge 29$M)]
\[
T = \frac{27{,}78 - 29{,}00}{3{,}76/\sqrt{40}} = \frac{-1{,}22}{0{,}595} \approx -2{,}05
\]
Valor crítico (cola izquierda, $\alpha=0{,}05$, $df=39$): $t_{0{,}05,\,39} \approx -1{,}685$.

Como $T = -2{,}05 < -1{,}685$, \textbf{se rechaza $H_0$} a nivel $\alpha = 0{,}05$. Los datos respaldan la afirmación del Director de Logística: el costo promedio real parece ser menor a 29 millones. El $p$-valor aproximado es 0{,}024 (menor que 0{,}05, consistente con el rechazo).
\end{resultadobox}

\begin{resultadobox}[Pregunta 3: Tamaño de Muestra]
Con margen de error $E = 0{,}5$ M, confianza del 95\% ($z_{0{,}975} = 1{,}96$) y $\sigma \approx s = 3{,}76$ M:
\[
n \approx \left(\frac{1{,}96 \times 3{,}76}{0{,}5}\right)^2 = \left(\frac{7{,}37}{0{,}5}\right)^2 = (14{,}74)^2 \approx 217{,}3 \implies n_{\min} = \mathbf{218} \text{ envíos}
\]
Monitorear 218 envíos garantiza una precisión de $\pm 0{,}5$ M COP con 95\% de confianza, asumiendo que la variabilidad futura es similar a la observada ($s \approx 3{,}76$ M).
\end{resultadobox}

\section*{Referencias}
\begin{enumerate}[label={[\arabic*]}]
    \item DeGroot, M.\,H. \& Schervish, M.\,J. (2012). \textit{Probability and Statistics}, 4th ed. Pearson. Cap.~8.
    \item Casella, G. \& Berger, R.\,L. (2002). \textit{Statistical Inference}, 2nd ed. Duxbury. Cap.~9.
    \item Law, A.\,M. (2015). \textit{Simulation Modeling and Analysis}, 5th ed. McGraw-Hill. Cap.~9.
\end{enumerate}

\end{document}
